% Context from chapters-tex/sparse-theory.tex; for mathematical reference, not compilation.
\section{Existence and Uniqueness}


                                                                                  
\subsection{Existence}

We study the penalized problem~\eqref{eq-lasso-lagr-ip}, with $\lambda>0$, and its constrained counterpart~\eqref{eq-lasso-constr-ip}:
\eql{\label{eq-lasso-lagr}\tag{$\Pp_\la(y)$}
	\umin{x \in \RR^N} f_\la(x) \eqdef \frac{1}{2\la} \norm{y-Ax}^2 + \norm{x}_1
} 
and the limiting problem as $\la \rightarrow 0$,
\eql{\label{eq-lasso-constr}\tag{$\Pp_0(y)$}
	\umin{A x = y} \norm{x}_1
	= \umin{x} f_0(x) \eqdef \iota_{\Ll_y}(x) + \norm{x}_1.
} 
where $A \in \RR^{P \times N}$, and $\Ll_y \eqdef \enscond{x \in \RR^N}{Ax=y}$.

Recall that we observe noisy measurements
\eq{
	y=A x_0 + w
}
We seek conditions ensuring that $x_0$ solves $\Pp_0(Ax_0)$ in the noiseless case, and quantitative bounds on its distance to solutions of $\Pp_\lambda(Ax_0+w)$ when $\lambda$ is chosen according to the noise level.

The objective in~\eqref{eq-lasso-lagr} is continuous and coercive because $\norm{\cdot}_1$ is coercive, so a minimizer exists.
 
Uniqueness is not automatic: $A$ may have a nontrivial kernel, and $\norm{\cdot}_1$ is not strictly convex.
 
If $y \in \Im(A)$, the constraint set in~\eqref{eq-lasso-constr} is nonempty. A constrained minimizer also exists but need not be unique.



\begin{figure}
\centering
  
\begin{tabular}{@{}c@{}c@{}}
\includegraphics[width=.35\linewidth]{sparse-theory/homotopy/homotopy-s3}&
\includegraphics[width=.35\linewidth]{sparse-theory/homotopy/homotopy-s6}\\
$s=3$  & $s=6$ \\
\includegraphics[width=.35\linewidth]{sparse-theory/homotopy/homotopy-s13}&
\includegraphics[width=.35\linewidth]{sparse-theory/homotopy/homotopy-s25}\\
$s=13$  & $s=25$ 
\end{tabular}
  
\caption{\label{fig-homotopy}
Solution paths $\la \mapsto x_\la$ for~\eqref{eq-lasso-lagr}.
}
\end{figure}

Figure~\ref{fig-homotopy} shows how the recovered coefficients change with $\lambda$, illustrating the effect of regularization on sparsity and reconstruction accuracy.





                                                                                  
\subsection{Polytope Projection for the Constrained Problem}
\label{sec-polytope-proj}

The next proposition characterizes noiseless $\ell^1$ recovery geometrically.

\begin{prop}\label{prop-polytope-proj}
	We denote $B \eqdef \enscond{x \in \RR^N}{\norm{x}_1 \leq 1}$. For $x_0\neq0$, assuming $\Im(A)=\RR^P$,
	\eql{\label{eq-polytop-ident} 
		x_0 \text{ is a solution to } \Pp_0(Ax_0)
		\quad\Longleftrightarrow\quad
		A \frac{x_0}{\norm{x_0}_1} \in \partial (AB)
	} 
	where ``$\partial$'' denotes the boundary and $AB=\enscond{Ax}{x\in B}$.
\end{prop}

\begin{proof}
Set $r=\norm{x_0}_1>0$ and $q=Ax_0/r$. If $x_0$ is not optimal, there is $z$ with $Az=Ax_0$ and $\norm{z}_1<r$. Since $A$ is onto, $A^+$ is a bounded right inverse. For sufficiently small $h\in\RR^P$,
\eq{
 \norm{z/r+A^+h}_1\leq\norm{z}_1/r+\norm{A^+h}_1<1,
}
and $q+h=A(z/r+A^+h)\in AB$. Thus $q$ is an interior point of $AB$.

Conversely, if $q\neq0$ is interior, then $(1+t)q\in AB$ for some $t>0$. Choose $u\in B$ with $Au=(1+t)q$. The vector $z=ru/(1+t)$ satisfies $Az=Ax_0$ and $\norm{z}_1<r$, so $x_0$ is not optimal. If $q=0$, the zero vector is already a better feasible point. These two implications prove the equivalence.
\begin{figure}[tbp]
\centering
\BookDrawing{tikz/sparse-theory-polytope-proof}
\caption{Geometry of the polytope characterization of exact recovery.}
\label{fig-polytope-proof}
\end{figure}
\end{proof}

For a nonzero recovered vector $x_0$, its measurement $Ax_0$ lies on the boundary of the scaled polytope $\norm{x_0}_1 A B$. When $P$ is much smaller than $N$, projection can remove faces and prevent $\ell^1$ recovery. Chapter~\ref{chap-cs} studies random projections that preserve enough of the geometry of the $\ell^1$ ball in $\RR^N